% !TEX encoding = UTF-8
\documentclass[12pt,a4paper]{article}

% ── Codificação e idioma ──
\usepackage[utf8]{inputenc}
\usepackage[T1]{fontenc}
\usepackage[brazil]{babel}

% ── Tipografia e layout ──
\usepackage{lmodern}
\usepackage[top=2.5cm, bottom=2.5cm, left=2.5cm, right=2.5cm]{geometry}
\usepackage{setspace}
\onehalfspacing
\usepackage{parskip}

% ── Gráficos ──
\usepackage{pgfplots}
\pgfplotsset{compat=1.18}
\usepgfplotslibrary{groupplots}
\usepackage{tikz}

% ── Tabelas ──
\usepackage{booktabs}
\usepackage{array}
\usepackage{longtable}

% ── Matemática ──
\usepackage{amsmath}

% ── Links e referências ──
\usepackage[colorlinks=true, linkcolor=blue, urlcolor=blue, citecolor=blue]{hyperref}

% ── Cabeçalhos ──
\usepackage{fancyhdr}
\pagestyle{fancy}
\fancyhf{}
\fancyhead[L]{\small Preditor de Preços de Imóveis}
\fancyhead[R]{\small Introdução ao Aprendizado de Máquina}
\fancyfoot[C]{\thepage}
\renewcommand{\headrulewidth}{0.4pt}

% ── Código inline ──
\usepackage{listings}
\lstset{
  basicstyle=\ttfamily\small,
  breaklines=true,
  frame=none
}

% ══════════════════════════════════════════════════════════════
\begin{document}

% ── Capa ──
\begin{titlepage}
\centering
\vspace*{3cm}

{\Large\bfseries Introdução ao Aprendizado de Máquina\par}
\vspace{1.5cm}
{\LARGE\bfseries Preditor de Preços de Imóveis\par}
\vspace{1cm}
{\large Projeto e implementação de um regressor de preços de imóveis\\
com pipeline de pré-processamento, limpeza orientada por dados,\\
comparação de modelos e arquitetura de roteamento\\
entre generalista e especialista\par}
\vspace{2cm}
{\large Aluno: Henrique Kezen V. H. Leite\par}
\vspace{0.5cm}
{Repositório:\par}
\vspace{0.5cm}
{Ambiente técnico: Python 3.14 --- Pandas, NumPy, Scikit-learn,\\
LightGBM, XGBoost, CatBoost, Matplotlib, Seaborn\par}
\vfill
{Rio de Janeiro, julho de 2026\par}
\end{titlepage}

% ── Sumário ──
\tableofcontents
\newpage

% ══════════════════════════════════════════════════════════════
\section*{Resumo}
\addcontentsline{toc}{section}{Resumo}

Este trabalho aborda a previsão de preços de imóveis a partir de um conjunto de 4.683 registros com 20 atributos. Foram testados Regressão Linear, XGBoost, LightGBM e CatBoost, todos treinados com alvo em \lstinline{log1p(preco)} --- cujas previsões são convertidas de volta para reais com \lstinline{expm1} --- variáveis categóricas nativas e \lstinline{sample_weight = 1/preco^1,5}. O modelo final é uma arquitetura de roteamento que combina um generalista (blend de 60\% CatBoost, 20\% XGBoost e 20\% LightGBM) com um especialista CatBoost para imóveis baratos (\lstinline{sample_weight = 1/preco²}), selecionados por um juiz contínuo.

% ══════════════════════════════════════════════════════════════
\section{Objetivo do Projeto}

O objetivo deste trabalho é construir um modelo de regressão capaz de prever o preço de venda de imóveis a partir de suas características físicas, localização e diferenciais de infraestrutura. A métrica de avaliação da competição é o \textbf{RMSPE} (\textit{Root Mean Squared Percentage Error}) --- a raiz da média dos erros percentuais ao quadrado. O RMSPE penaliza proporcionalmente ao preço real: um erro de R\$\,50 mil em um imóvel de R\$\,100 mil é muito mais grave que o mesmo erro em um imóvel de R\$\,1 milhão.

% ══════════════════════════════════════════════════════════════
\section{Descrição dos Dados}

O conjunto de dados fornecido pela disciplina está dividido em duas partes:

\begin{table}[h]
\centering
\begin{tabular}{lrcl}
\toprule
\textbf{Arquivo} & \textbf{Registros} & \textbf{Colunas} & \textbf{Uso} \\
\midrule
\texttt{conjunto\_de\_treinamento.csv} & 4.683 & 21 & Treino e validação \\
\texttt{conjunto\_de\_teste.csv}       & 2.001 & 20 & Previsões para o Kaggle \\
\bottomrule
\end{tabular}
\end{table}

\subsection{Variável alvo}

A variável \texttt{preco} é contínua e apresenta distribuição fortemente assimétrica à direita: a média (R\$\,927.705) é quase o dobro da mediana (R\$\,515.000), com mínimo de R\$\,750 e máximo de R\$\,630 milhões. Essa assimetria motivou o uso da transformação logarítmica \lstinline{log1p(preco)} como alvo de treino.

\subsection{Tipos de variáveis}

\begin{table}[h]
\centering
\begin{tabular}{lrl}
\toprule
\textbf{Tipo} & \textbf{Qtd} & \textbf{Exemplos} \\
\midrule
Categóricas nominais     & 4  & \texttt{tipo}, \texttt{bairro}, \texttt{tipo\_vendedor}, \texttt{diferenciais} \\
Numéricas quantitativas   & 5  & \texttt{quartos}, \texttt{suites}, \texttt{vagas}, \texttt{area\_util}, \texttt{area\_extra} \\
Binárias (diferenciais)   & 10 & \texttt{churrasqueira}, \texttt{piscina}, \texttt{vista\_mar}, \texttt{sauna} \\
Identificador             & 1  & \texttt{Id} (não preditivo) \\
\bottomrule
\end{tabular}
\end{table}

A coluna \texttt{tipo} é altamente desbalanceada: 4.501 apartamentos contra apenas 177 casas, 3 lofts e 2 quitinetes. A coluna \texttt{bairro} possui 66 categorias distintas, das quais 30 possuem menos de 10 imóveis.

\subsection{Valores faltantes}

Nenhuma coluna apresentou valores nulos, o que simplificou o pré-processamento.

\subsection{Valores extremos}

Variáveis numéricas apresentaram outliers significativos que exigiram análise caso a caso:

\begin{table}[h]
\centering
\begin{tabular}{lrrl}
\toprule
\textbf{Variável} & \textbf{Mediana} & \textbf{Máximo} & \textbf{Observação} \\
\midrule
\texttt{preco}       & R\$\,515.000 & R\$\,630.000.000 & Provável erro de escala \\
\texttt{area\_util}  & $\sim$100\,m² & 2.045\,m²        & Casas decimais perdidas \\
\texttt{area\_extra} & 0\,m²         & 17.450\,m²       & Provável erro de escala \\
\texttt{vagas}       & 2             & 30                & Caso isolado \\
\bottomrule
\end{tabular}
\end{table}

% ══════════════════════════════════════════════════════════════
\section{Metodologia}

\subsection{Estratégia geral}

O projeto foi organizado em fases: pré-processamento e auditoria dos dados, modelagem com comparação de algoritmos, otimização de hiperparâmetros, construção de um modelo especialista para imóveis baratos e, por fim, arquitetura de roteamento com juiz. O conjunto de teste foi mantido isolado até a etapa final de submissão.

\subsection{Métricas de validação}

Ao longo do relatório, três grandezas diferentes são reportadas:

\begin{itemize}
  \item \textbf{Validação interna (KFold)} --- RMSPE médio da validação cruzada de 5 folds sobre os dados de treino, com \lstinline{shuffle=True} e \lstinline{random_state=42}. É a métrica usada para tomar decisões de modelagem.
  \item \textbf{Holdout por blocos de ID} --- RMSPE calculado sobre os 20\% menores IDs do treino, simulando a direção treino$\to$teste observada na competição (IDs de teste vão de 0 a $\sim$1.999, IDs de treino de $\sim$2.000 a $\sim$6.682). Esse holdout complementa o KFold porque a separação do Kaggle não é aleatória.
  \item \textbf{Score público (Kaggle)} --- RMSPE calculado pela plataforma sobre a fração pública do conjunto de teste. Serve como estimativa externa de generalização.
\end{itemize}

\subsection{Comparação de modelos}

\begin{itemize}
  \item \textbf{Regressão Linear} --- baseline de referência, estabelecendo um piso de desempenho.
  \item \textbf{XGBoost} --- gradient boosting com crescimento por profundidade.
  \item \textbf{LightGBM} --- gradient boosting com crescimento por folha, usando \texttt{bairro} como categoria nativa.
  \item \textbf{CatBoost} --- gradient boosting com ordered boosting e tratamento nativo de todas as categorias, sem necessidade de one-hot encoding.
\end{itemize}

\subsection{Métrica RMSPE}

\begin{equation}
\text{RMSPE} = \sqrt{\frac{1}{n}\sum_{i=1}^{n}\left(\frac{y_i - \hat{y}_i}{y_i}\right)^2}
\end{equation}

Como os erros são divididos pelo preço real, imóveis baratos com o mesmo erro absoluto produzem erro percentual muito maior. Essa propriedade motivou o uso de sample weights ao longo do projeto.

% ══════════════════════════════════════════════════════════════
\section{Pré-processamento dos Dados}

\subsection{Objetivo}

Transformar a base bruta em uma matriz numérica pronta para modelagem, tratando outliers, criando features derivadas e codificando variáveis categóricas.

\subsection{Auditoria e limpeza de outliers}

A auditoria foi feita em duas fases. Primeiro, inspeção direta de preços e áreas extremas. Segundo, criação de \textbf{features de auditoria} que utilizam o preço --- como \lstinline{preco_m2 = preco / area_util} --- para identificar inconsistências menos óbvias. Essas features existem exclusivamente para análise exploratória e \textbf{não são entregues ao modelo}, já que o preço é justamente a variável que o modelo deve prever e, portanto, não está disponível no conjunto de teste. Os casos individuais foram comparados com a mediana de preço/m² do respectivo bairro.

\subsubsection{Primeira limpeza e o erro de ser agressivo demais}

Na primeira versão, 29 imóveis foram removidos: 6 casos de erro claro (preços absurdos, áreas impossíveis) e 23 casos cujo preço por metro quadrado destoava da mediana do bairro. Essa remoção parecia segura porque a validação cruzada aleatória (KFold) não piorava.

Entretanto, ao submeter ao Kaggle, o resultado foi pior do que o esperado. A investigação revelou que \textbf{parte dos imóveis removidos possuía registros com features muito similares no conjunto de teste} --- ou seja, o teste continha registros praticamente iguais a esses casos que havíamos descartado. Sem esses exemplos no treino, o modelo não tinha como prever corretamente esses imóveis repetidos.

Além disso, o KFold aleatório não capturava essa divergência porque a separação treino/teste no Kaggle não é aleatória: os IDs de teste são sistematicamente menores ($\sim$0 a $\sim$1.999) que os de treino ($\sim$2.000 a $\sim$6.682). Essa descoberta levou à adoção do \textbf{holdout por blocos de ID} como validação complementar.

\subsubsection{Recuperação e correção conservadora}

As 23 linhas foram então recuperadas. Ao invés de simplesmente recolocá-las sem tratamento, investigou-se a causa dos valores extremos. Descobriu-se que 21 delas eram apartamentos cuja \texttt{area\_util} tinha provável erro de casa decimal --- a regra \lstinline{area_util / quartos > 200} os identificava deterministicamente, sem nenhum outro apartamento do treino sendo selecionado. A correção consistiu em dividir \texttt{area\_util} por 10.

Adicionalmente:

\begin{itemize}
  \item Dois preços foram corrigidos por divisão por 10 (IDs 2749 e 4316), com base na comparação com a mediana do bairro;
  \item O número de vagas do ID 6383 foi corrigido de 30 para 3;
  \item A área extra do ID 6654 foi corrigida dividindo por 100, tornando o registro plausível e permitindo sua reinclusão.
\end{itemize}

\subsubsection{Exclusões definitivas e resultado}

Cinco imóveis permaneceram excluídos por não possuírem correção confiável:

\begin{table}[h]
\centering
\begin{tabular}{rl}
\toprule
\textbf{ID} & \textbf{Motivo da exclusão definitiva} \\
\midrule
5910 & Preço de R\$\,750 \\
2405 & Preço de R\$\,630 milhões para 98\,m² \\
4568 & Preço de R\$\,65 milhões para 36\,m² \\
6004 & Preço de R\$\,340 milhões para 72\,m² \\
6654 & 17.450\,m² de área extra \\
\bottomrule
\end{tabular}
\end{table}

A mesma regra de \lstinline{area_util / quartos > 200} foi aplicada ao conjunto de teste, corrigindo 3 imóveis. No conjunto de teste, 433 dos 2.000 imóveis possuem combinação completa de features muito similar a pelo menos um do treino --- incluindo vários dos casos anteriormente removidos. Isso confirmou que a remoção agressiva inicial prejudicava diretamente as previsões desses imóveis repetidos.

A base final ficou com \textbf{4.678 imóveis} para treino.

\subsection{Transformação do alvo}

O preço foi transformado com \lstinline{np.log1p(preco)} para o treinamento: o modelo aprende a prever nessa escala comprimida. Depois da previsão, o valor é convertido de volta para reais com \lstinline{np.expm1}, recuperando a escala original do preço antes de qualquer avaliação ou submissão. Essa transformação é especialmente eficaz para a métrica RMSPE porque converte divisão em subtração: minimizar o MSE no espaço log equivale aproximadamente a minimizar erros proporcionais.

\subsection{Variáveis categóricas}

\begin{table}[h]
\centering
\begin{tabular}{lll}
\toprule
\textbf{Coluna} & \textbf{Tratamento} & \textbf{Justificativa} \\
\midrule
\texttt{tipo}           & One-hot (Loft + Quitinete $\to$ Outro) & Apenas 3 e 2 exemplos \\
\texttt{bairro}         & Categoria nativa; $<$10 imóveis $\to$ ``Outros'' & 91 $\to$ 26 features \\
\texttt{tipo\_vendedor} & Binário: Imobiliária=1 & Apenas duas categorias \\
\texttt{diferenciais}   & Removida & Redundante com binárias \\
\bottomrule
\end{tabular}
\end{table}

\subsection{Feature engineering}

\begin{table}[h]
\centering
\begin{tabular}{lll}
\toprule
\textbf{Feature} & \textbf{Fórmula} & \textbf{Motivação} \\
\midrule
\texttt{area\_total}       & \texttt{area\_util + area\_extra} & Tamanho total \\
\texttt{area\_por\_quarto} & \texttt{area\_util / quartos}     & Amplitude dos cômodos \\
\texttt{n\_comodidades}    & Soma das 10 binárias              & Proxy de luxo \\
\texttt{tem\_suite}        & 1 se \texttt{suites > 0}          & Diferencial de qualidade \\
\texttt{vagas\_por\_quarto}& \texttt{vagas / quartos}           & Conveniência relativa \\
\bottomrule
\end{tabular}
\end{table}

\subsection{Resultado}

A matriz processada final para o LightGBM/CatBoost categórico ficou com 4.678 linhas $\times$ 26 features, inteiramente numérica e pronta para modelagem.

% ══════════════════════════════════════════════════════════════
\section{Modelagem e Treinamento}

\subsection{Objetivo}

Comparar o desempenho dos modelos de regressão e identificar os mais promissores para refinamento e composição.

\subsection{Resultados comparativos iniciais}

Os modelos foram avaliados sobre os mesmos 931 imóveis de um holdout fixo (80/20, \lstinline{random_state=42}). A Figura~\ref{fig:comparacao_modelos} apresenta o RMSPE de cada modelo.

\begin{figure}[h]
\centering
\begin{tikzpicture}
\begin{axis}[
    xbar,
    width=12cm,
    height=6cm,
    xlabel={RMSPE (\%)},
    symbolic y coords={
        {LightGBM (log)},
        {XGBoost (log)},
        {Reg. Linear (log)},
        {Reg. Linear (normal)}
    },
    ytick=data,
    xmin=20, xmax=45,
    nodes near coords,
    nodes near coords align={horizontal},
    nodes near coords style={font=\small, /pgf/number format/.cd, fixed, precision=2},
    bar width=12pt,
    enlarge y limits=0.25,
    every axis plot/.append style={fill=blue!60},
]
\addplot coordinates {
    (40.21,{Reg. Linear (normal)})
    (30.13,{Reg. Linear (log)})
    (25.26,{XGBoost (log)})
    (25.05,{LightGBM (log)})
};
\end{axis}
\end{tikzpicture}
\caption{Comparação do RMSPE dos quatro modelos no holdout fixo. A transformação logarítmica reduziu o erro em mais de 10 pontos percentuais na regressão linear. Os modelos de boosting superaram a regressão em log por aproximadamente 5 pontos.}
\label{fig:comparacao_modelos}
\end{figure}

A transformação logarítmica reduziu o RMSPE em 10,08 pontos percentuais e eliminou as 10 previsões negativas da regressão normal. Os modelos de boosting superaram a regressão linear em aproximadamente 5 pontos, mas a diferença entre XGBoost e LightGBM foi de apenas 0,21 ponto --- dentro do ruído da amostra.

\subsection{Impacto da representação do bairro}

A mudança de one-hot encoding (91 features) para categoria nativa (26 features) trouxe ganho consistente, conforme ilustrado na Figura~\ref{fig:bairro_categorico}.

\begin{figure}[h]
\centering
\begin{tikzpicture}
\begin{axis}[
    ybar,
    width=9cm,
    height=6cm,
    ylabel={RMSPE médio (\%)},
    symbolic x coords={One-hot (91 feat.), Categórico (26 feat.)},
    xtick=data,
    ymin=22.5, ymax=25,
    nodes near coords,
    nodes near coords style={font=\small, /pgf/number format/.cd, fixed, precision=2},
    bar width=25pt,
    enlarge x limits=0.4,
    every axis plot/.append style={fill=blue!60},
    error bars/y dir=both,
    error bars/y explicit,
]
\addplot+[fill=red!40] coordinates {
    (One-hot (91 feat.), 24.08) +- (0, 1.51)
};
\addplot+[fill=green!50] coordinates {
    (Categórico (26 feat.), 23.39) +- (0, 1.37)
};
\legend{One-hot, Categórico}
\end{axis}
\end{tikzpicture}
\caption{RMSPE médio com barras de desvio-padrão ($\pm$1\,d.p. entre folds). A representação categórica nativa superou o one-hot em todos os cinco folds, com ganhos entre 0,06 e 1,61 ponto percentual.}
\label{fig:bairro_categorico}
\end{figure}

\subsection{Otimização de hiperparâmetros}

O LightGBM foi otimizado com \lstinline{RandomizedSearchCV} (40 combinações $\times$ 5 folds = 200 treinamentos). A melhor combinação de taxa de aprendizado e número de árvores foi \lstinline{learning_rate=0,03} com 600--800 árvores, atingindo RMSPE médio de 23,34\%.

% ══════════════════════════════════════════════════════════════
\section{Correspondências estruturais e pós-processamento}

\subsection{Repetições entre treino e teste}

Conforme descrito na auditoria de dados (Seção 4.2), 433 dos 2.000 imóveis de teste possuem combinação completa de features muito similar a pelo menos um imóvel do treino, formando 843 pares praticamente iguais. Isso indica repetição de anúncios ou registros no dataset. As chaves de correspondência utilizadas são:

\begin{center}
\texttt{tipo, bairro, area\_util, area\_extra, quartos, suites, vagas}
\end{center}

\subsection{Pós-processamento por correspondência}

Para imóveis com correspondência estrutural, a previsão final combina a saída do modelo com a mediana de preço dos registros correspondentes no treino:

\begin{equation}
\hat{y}_{\text{final}} = (1 - w) \cdot \hat{y}_{\text{modelo}} + w \cdot \text{mediana}_{\text{correspondência}}
\end{equation}

A mediana de correspondência é calculada somente com os dados de treino de cada fold, evitando vazamento do alvo. Na base final, 988 dos 2.000 imóveis de teste ($\sim$49\%) encontraram pelo menos uma correspondência. O peso $w = 0{,}30$ foi o melhor internamente e confirmou-se no Kaggle, produzindo um platô estável entre 25\% e 40\%.

% ══════════════════════════════════════════════════════════════
\section{Combinação de modelos (Blend) e pesos por preço}

\subsection{Conceito do blend}

Como XGBoost e LightGBM obtiveram resultados muito próximos na modelagem inicial, foi testada a combinação de suas previsões por média simples --- uma técnica chamada \textbf{blend}. A ideia é que, por usarem estratégias diferentes de construção de árvores e tratamento de categorias, os dois modelos tendem a errar em imóveis diferentes. Ao calcular a média das previsões, os erros individuais se compensam parcialmente, reduzindo a variância sem aumentar o viés.

O blend de 50\% XGBoost e 50\% LightGBM obteve \textbf{23,14\%} no KFold, melhor que ambos os modelos isolados. Esse resultado motivou o uso de blends em todos os experimentos subsequentes.

\subsection{Sample weights: alinhando treino e métrica}

Como o RMSPE penaliza erros percentuais, imóveis baratos dominam a métrica. A função de perda padrão (MSE sobre log) já aproxima erros proporcionais, mas não dá prioridade explícita a nenhuma faixa de preço. A introdução de $\text{sample\_weight} = 1 / \text{preco}^\alpha$, normalizado para média 1, força o modelo a ``se esforçar mais'' nos imóveis baratos.

\subsection{Resultados do blend XGBoost/LightGBM com diferentes alphas}

A Figura~\ref{fig:sample_weights} mostra a evolução do RMSPE em função do expoente $\alpha$ dos sample weights, tanto no KFold quanto no holdout por blocos de ID.

\begin{figure}[h]
\centering
\begin{tikzpicture}
\begin{axis}[
    width=12cm,
    height=7cm,
    xlabel={$\alpha$ (expoente do sample weight)},
    ylabel={RMSPE (\%)},
    xmin=-0.1, xmax=1.7,
    ymin=20.5, ymax=23.5,
    xtick={0, 0.5, 1.0, 1.5},
    legend pos=north east,
    grid=major,
    grid style={dashed, gray!30},
    mark size=3pt,
    thick,
]
\addplot[color=blue, mark=square*, thick] coordinates {
    (0.00, 23.14)
    (0.50, 22.32)
    (1.00, 21.73)
    (1.50, 21.54)
};
\addlegendentry{KFold (5 folds)}

\addplot[color=red, mark=triangle*, thick] coordinates {
    (0.00, 21.86)
    (0.50, 21.42)
    (1.00, 20.96)
    (1.50, 20.91)
};
\addlegendentry{Holdout por ID}

\end{axis}
\end{tikzpicture}
\caption{RMSPE do blend XGBoost/LightGBM em função de $\alpha$. O ganho veio principalmente dos imóveis baratos: no quartil até R\$\,355 mil, o RMSPE caiu de 27,34\% ($\alpha=0$) para 23,78\% ($\alpha=1{,}5$). Em contrapartida, o quartil mais caro subiu de 22,62\% para 23,77\%.}
\label{fig:sample_weights}
\end{figure}

No Kaggle, o blend com $\alpha=1{,}5$ obteve \textbf{21,86\%}, contra 22,63\% com $\alpha=1$ --- confirmando externamente a vantagem dos pesos mais agressivos.

% ══════════════════════════════════════════════════════════════
\section{CatBoost e Blend Global}

\subsection{CatBoost como terceiro modelo}

O CatBoost foi testado com categorias totalmente nativas (\texttt{tipo}, \texttt{bairro}, \texttt{tipo\_vendedor}, \texttt{diferenciais}), dispensando qualquer codificação prévia. Com \lstinline{sample_weight = 1/preco} e correspondência de 30\%, obteve RMSPE de 21,97\% no KFold --- competitivo com o blend LightGBM/XGBoost.

No Kaggle, o CatBoost com $\alpha=1{,}5$ e correspondência de 30\% obteve \textbf{21,82\%}.

\subsection{Blend global de três modelos}

A combinação de 60\% CatBoost + 20\% XGBoost + 20\% LightGBM, todos com $\alpha=1{,}5$ e correspondência de 30\%, obteve \textbf{21,45\%} no Kaggle --- melhorando 0,37 ponto sobre o CatBoost isolado e 0,41 ponto sobre o blend XGBoost/LightGBM. Esse resultado confirmou que os modelos cometem erros complementares e estabeleceu a referência \textbf{generalista}.

% ══════════════════════════════════════════════════════════════
\section{Especialista para imóveis baratos}

\subsection{Motivação}

Apesar da melhora trazida pelos sample weights, o generalista ainda apresentava RMSPE alto no quartil de imóveis baratos (até R\$\,355 mil). Um CatBoost dedicado, com ponderação extrema (\lstinline{sample_weight = 1/preco²}), foi buscado para essa faixa.

\subsection{Busca e configuração final}

Foram realizados 85 treinamentos em sete blocos experimentais. O melhor especialista:

\begin{table}[h]
\centering
\begin{tabular}{ll}
\toprule
\textbf{Parâmetro} & \textbf{Valor} \\
\midrule
Alvo                  & \lstinline{log1p(preco)} \\
\texttt{sample\_weight} & $1 / \text{preco}^2$ \\
\texttt{depth}         & 7 \\
\texttt{iterations}    & 600 \\
\texttt{random\_strength} & 0 \\
Bootstrap              & Bayesiano \\
Correspondência        & 0\% \\
Fator de calibração    & 0,914489 \\
\bottomrule
\end{tabular}
\end{table}

\subsection{Resultado do especialista}

\begin{figure}[h]
\centering
\begin{tikzpicture}
\begin{axis}[
    ybar,
    width=9cm,
    height=5.5cm,
    ylabel={RMSPE até R\$\,355 mil (\%)},
    symbolic x coords={Cat original ($\alpha$=1.5), Especialista final},
    xtick=data,
    ymin=18, ymax=20,
    nodes near coords,
    nodes near coords style={font=\small, /pgf/number format/.cd, fixed, precision=2},
    bar width=25pt,
    enlarge x limits=0.4,
]
\addplot+[fill=red!40] coordinates {(Cat original ($\alpha$=1.5), 19.29)};
\addplot+[fill=green!50] coordinates {(Especialista final, 18.90)};
\end{axis}
\end{tikzpicture}
\caption{RMSPE calibrado cross-fit no quartil de imóveis até R\$\,355 mil. A melhora apareceu em quatro dos cinco folds, com intervalo indicativo de 0,14 a 0,66 ponto percentual.}
\label{fig:especialista}
\end{figure}

% ══════════════════════════════════════════════════════════════
\section{Arquitetura de roteamento com juiz}

\subsection{Conceito}

A ideia é usar o generalista para a maioria dos imóveis e substituí-lo parcialmente pelo especialista nos imóveis onde esse último é superior. Um modelo ``juiz'' aprende a estimar, para cada linha, o peso ótimo do especialista.

\subsection{Juiz binário (classificador de utilidade)}

Um \lstinline{CatBoostClassifier} aprendeu se o especialista produziu menor perda percentual quadrática que o generalista, ponderado pelo impacto absoluto da decisão. A previsão final é:

\begin{equation}
\hat{y} = G + \lambda \cdot w_{\text{juiz}} \cdot (S - G)
\end{equation}

\noindent onde $G$ é a previsão do generalista, $S$ a do especialista, $w_{\text{juiz}}$ o peso previsto pelo juiz e $\lambda$ a intensidade do roteamento.

Quatro juízes supervisionados foram comparados em cross-fit. O classificador de utilidade venceu com RMSPE de 21,3394\%, contra 21,4066\% do generalista sem juiz. No Kaggle, com intensidade $\lambda=1{,}00$, obteve \textbf{21,33\%}.

\subsection{Juiz contínuo}

O último refinamento substituiu a classificação binária por uma regressão que estima diretamente o peso ótimo do especialista, usando como pseudo-alvo $(y-G)/D$ projetado em $[0,1]$ e ponderado por $(D/y)^2$.

\begin{table}[h]
\centering
\begin{tabular}{lrr}
\toprule
\textbf{Modelo} & \textbf{RMSPE cross-fit} & \textbf{RMSPE holdout por ID} \\
\midrule
Juiz binário (melhor intensidade) & 21,3394\% & 21,0035\% \\
\textbf{Juiz contínuo calibrado}  & \textbf{21,2873\%} & \textbf{20,9536\%} \\
\bottomrule
\end{tabular}
\end{table}

O juiz contínuo superou o binário nos cinco folds, com ganhos entre 0,012 e 0,074 ponto percentual.

% ══════════════════════════════════════════════════════════════
\section{Discussão dos Resultados}

\subsection{Por que a transformação logarítmica foi essencial?}

A transformação \lstinline{log1p} reduziu o RMSPE em mais de 10 pontos percentuais na regressão linear e eliminou previsões negativas. A explicação é que o RMSPE avalia erro relativo, e o log transforma divisão em subtração --- ao minimizar o MSE no espaço log, o modelo aprende proporções naturalmente. Todos os modelos subsequentes foram treinados com esse alvo.

\subsection{Por que a categoria nativa superou o one-hot?}

O one-hot encoding criava 66 colunas binárias para os bairros, dificultando que as árvores tratassem em conjunto bairros com comportamento semelhante e favorecendo a memorização de bairros com poucos exemplos. A representação categórica nativa permite ao LightGBM/CatBoost procurar separações ótimas entre grupos de bairros durante a construção de cada árvore, obtendo mais informação com menos dimensões.

\subsection{Por que sample weights ajudaram tanto?}

O RMSPE é uma métrica assimétrica em relação ao preço: um erro de R\$\,50 mil em um imóvel de R\$\,100 mil contribui 25 vezes mais que o mesmo erro em um imóvel de R\$\,500 mil. A função de perda padrão (MSE no log) já aproxima erros proporcionais, mas trata todas as faixas igualmente. Os sample weights $1/\text{preco}^\alpha$ fazem o modelo gastar mais capacidade nos imóveis baratos, alinhando explicitamente treino e avaliação. O ganho foi de mais de 2 pontos percentuais no Kaggle.

\subsection{Por que o blend global foi superior aos modelos isolados?}

O CatBoost, XGBoost e LightGBM usam estratégias diferentes de construção de árvores, tratamento de categorias e regularização. Isso faz com que cometam erros em imóveis diferentes. A média ponderada das previsões reduziu a variância sem aumentar o viés, melhorando o RMSPE final.

\subsection{Por que o juiz contínuo superou o binário?}

O juiz binário decide apenas ``sim ou não'' para o uso do especialista, perdendo informação sobre o grau de confiança. O juiz contínuo estima diretamente o peso ótimo, permitindo transições suaves entre generalista e especialista. Isso é especialmente útil nos imóveis na fronteira entre ``baratos'' e ``intermediários'', onde nenhum dos dois modelos é claramente superior.

\subsection{Sobre a divergência entre validação interna e Kaggle}

A separação de treino e teste no Kaggle não é aleatória --- os IDs de teste são sistematicamente menores. Isso faz com que o KFold aleatório subestime ou superestime ganhos de certas estratégias (como a recuperação de linhas anômalas). A adoção do holdout por blocos de ID como validação complementar foi a medida mais importante para alinhar validação interna com o score público.

% ══════════════════════════════════════════════════════════════
\section{Evolução dos resultados no Kaggle}

A Figura~\ref{fig:evolucao_kaggle} apresenta a progressão cronológica do RMSPE público ao longo das principais submissões.

\begin{figure}[h]
\centering
\begin{tikzpicture}
\begin{axis}[
    xbar,
    width=14cm,
    height=10cm,
    xlabel={RMSPE no Kaggle (\%)},
    symbolic y coords={
        {Juiz binário},
        {Blend global 60/20/20},
        {CatBoost $\alpha$=1.5},
        {Blend XGB/LGB $\alpha$=1.5},
        {XGBoost target enc.},
        {LGB corrigido + corr. 30\%},
        {LGB 800 árv. categórico},
        {LGB ajustado},
        {LGB inicial}
    },
    ytick=data,
    xmin=20, xmax=29,
    nodes near coords,
    nodes near coords align={horizontal},
    nodes near coords style={font=\small, /pgf/number format/.cd, fixed, precision=2},
    bar width=10pt,
    enlarge y limits=0.08,
    y tick label style={font=\small},
    every axis plot/.append style={fill=blue!50},
]
\addplot coordinates {
    (27.85,{LGB inicial})
    (27.95,{LGB ajustado})
    (25.50,{LGB 800 árv. categórico})
    (24.76,{LGB corrigido + corr. 30\%})
    (24.86,{XGBoost target enc.})
    (21.86,{Blend XGB/LGB $\alpha$=1.5})
    (21.82,{CatBoost $\alpha$=1.5})
    (21.45,{Blend global 60/20/20})
    (21.33,{Juiz binário})
};
\end{axis}
\end{tikzpicture}
\caption{Evolução cronológica do RMSPE público no Kaggle. As maiores quedas vieram da introdução de sample weights ($\sim$3 p.p.) e da adição de dados corrigidos com correspondência estrutural ($\sim$2,6 p.p.).}
\label{fig:evolucao_kaggle}
\end{figure}

% ══════════════════════════════════════════════════════════════
\section{Conclusão}

O modelo final utilizado na competição é uma \textbf{arquitetura de roteamento} com a seguinte composição:

\begin{table}[h]
\centering
\begin{tabular}{ll}
\toprule
\textbf{Componente} & \textbf{Detalhes} \\
\midrule
\textbf{Generalista}   & Blend: 60\% CatBoost + 20\% XGBoost + 20\% LightGBM \\
Ponderação              & \lstinline{sample_weight = 1/preco^1.5} em todos os componentes \\
Correspondência         & 30\% da mediana estrutural para $\sim$49\% dos imóveis de teste \\
\textbf{Especialista}  & CatBoost depth 7, 600 iterações, $1/\text{preco}^2$, calibração 0,914 \\
\textbf{Juiz contínuo} & CatBoost regressão, alvo projetado $[0,1]$, calibrado cross-fit \\
Mapeamento final        & $\text{clip}(-0{,}45 + 1{,}50 \times \text{score\_juiz},\; 0,\; 1)$ \\
Alvo de treino          & \lstinline{log1p(preco)} \\
Dados de treino         & 4.678 imóveis (5 exclusões, correções conservadoras de escala) \\
Features                & 26 (com bairro como categoria nativa) \\
Melhor score Kaggle     & \textbf{21,33\%} \\
\bottomrule
\end{tabular}
\end{table}

As estratégias que se mostraram mais efetivas foram, em ordem de impacto:

\begin{enumerate}
  \item \textbf{Transformação logarítmica} do preço ($\sim$10\,p.p. de ganho na regressão linear);
  \item \textbf{Sample weights por preço} ($\sim$2\,p.p. de ganho no Kaggle);
  \item \textbf{Correspondência estrutural} de registros repetidos ($\sim$0,3--1,5\,p.p.);
  \item \textbf{Categoria nativa} para bairros ($\sim$0,7\,p.p.);
  \item \textbf{Blend de três modelos} ($\sim$0,4\,p.p.);
  \item \textbf{Roteamento com juiz} ($\sim$0,1\,p.p.);
  \item \textbf{Auditoria e correção conservadora} de erros de escala (melhora indireta via qualidade dos dados).
\end{enumerate}

Tentativas de ampliar a busca de hiperparâmetros do LightGBM ajustado, corrigir \texttt{area\_extra} em apartamentos e usar o XGBoost inicial sem otimização não trouxeram ganho no placar real.

% ══════════════════════════════════════════════════════════════
\section{Referências}

\begin{itemize}
  \item Scikit-learn: \url{https://scikit-learn.org/}
  \item LightGBM: \url{https://lightgbm.readthedocs.io/}
  \item XGBoost: \url{https://xgboost.readthedocs.io/}
  \item CatBoost: \url{https://catboost.ai/}
  \item Pandas: \url{https://pandas.pydata.org/}
\end{itemize}

\end{document}
